% !TeX encoding = UTF-8
\documentclass[variant=guia,cover=institutional,mode=screen]{../cls/ua-engineering-v6-article}
% !TeX encoding = UTF-8
\UASetUniversity{Universidad Austral}
\UASetFaculty{Facultad de Ingeniería}
\UASetDepartment{Departamento de Computación}
\UASetCampus{Pilar}
\UASetCareer{Ingeniería Informática}
\UASetCommission{Comisión A}
\UASetProfessor{Equipo docente}
\UASetSemester{Segundo cuatrimestre 2026}
\UASetCourse{Sistemas Distribuidos}
\UASetDate{2026}


\UASetClassNumber{02}
\UASetClassTitle{RPC, timeouts, retries e idempotencia}
\UASetDocKind{Guía resuelta}

\title{Clase 02 --- Guía resuelta}
\author{\UAProfessor}
\date{\UADate}

\begin{document}
\maketitle

\section{Respuestas seleccionadas}

\begin{uaexercise}
Definir RPC.
\end{uaexercise}

\begin{uaanswer}
RPC es un mecanismo que permite invocar procedimientos en otra máquina como si fueran locales, ocultando la comunicación de red.

El problema es que esta abstracción oculta latencia, fallas y duplicación, lo que puede llevar a errores de diseño.
\end{uaanswer}

\begin{uaexercise}
Construir una línea de tiempo donde el cliente abandona por timeout pero la operación remota igualmente puede ejecutarse.
\end{uaexercise}

\begin{uaanswer}
Una línea posible: $t_0$ el cliente envía request; $t_1$ el transporte retransmite o el cliente RPC reintenta; $t_2$ la aplicación vence su timeout y devuelve error; $t_3$ el servidor ejecuta; $t_4$ la base confirma; $t_5$ la respuesta llega tarde o se pierde. Después de $t_2$ todavía pueden ocurrir retransmisiones, ejecución, confirmación y respuesta tardía.
\end{uaanswer}

\begin{uaexercise}
Explicar por qué separar capas ayuda a diagnosticar una falla de comunicación.
\end{uaexercise}

\begin{uaanswer}
Porque cada capa observa una parte distinta del fenómeno y puede tomar acciones independientes. La aplicación puede ver timeout, el cliente HTTP puede haber reintentado, TCP puede haber retransmitido, el balanceador puede haber redirigido, el servidor puede haber ejecutado y la base puede haber confirmado. Separar capas evita atribuir todo a ``la red'' y permite identificar mitigaciones concretas.
\end{uaanswer}

\begin{uaexercise}
¿Por qué retry introduce duplicación?
\end{uaexercise}

\begin{uaanswer}
Porque el cliente no puede distinguir si la ejecución anterior ocurrió.

Formalmente:

\begin{itemize}
\item ejecución A: operación no ejecutada
\item ejecución B: operación ejecutada pero respuesta perdida
\end{itemize}

Ambas son indistinguibles → el cliente reintenta → posible doble ejecución.
\end{uaanswer}

\begin{uaexercise}
Diseñar idempotencia para pagos.
\end{uaexercise}

\begin{uaanswer}
Solución:

\begin{itemize}
\item generar request ID único
\item persistir resultado
\item rechazar duplicados
\end{itemize}

Esto convierte múltiples ejecuciones en una sola operación lógica.
\end{uaanswer}

\begin{uaexercise}
Aplicar el modelo D.
\end{uaexercise}

\begin{uaanswer}
\begin{itemize}
\item $P$: RPC
\item $F$: timeout + pérdida
\item $G$: evitar duplicación
\item $S$: idempotencia
\item $T$: más estado
\end{itemize}
\end{uaanswer}

\end{document}
